\scanpage{8}
\mspage{\handwritten{0}}

XII \quad \uline{Tores maximaux}, \uline{groupe de Weyl}, \uline{sous-groupes de Cartan}, \uline{centre réductif des schémas en groupes lisses et affines}.

par A. Grothendieck .

A partir du présent exposé, contrairement aux exposés précédents, nous ferons usage \deleted{\unsure{ce|le}}\handwritten{de}s résultats connus sur la structure des groupes algébriques lisses sur un corps algébriquement clos $k$, et surtout de la théorie de B\deleted{orel}\handwritten{OREL} des groupes algébriques affines, exposée dans le Séminaire Ch\deleted{evalley}\handwritten{EVALLEY} 56/57 : \deleted{sur les Groupes} ``Classification des groupes de Lie Algébriques'' . Suivant l'usage en théorie des groupes algébriques, nous référons à ce Séminaire par le sigle BIBLE. Pour les \deleted{moment,} \inserted{prochains exposés,} nous aurons besoin \inserted{notamment} des résultats de\deleted{s} BIBLE 4 , 5, 6, 7 (le chiffre après BIBLE renvoi\unsure{\handwritten{é}|e} au numéro de l'exposé).

Il \deleted{\unsure{me|ne}} semble \deleted{\unsure{pas}} \inserted{d'ailleurs} que la théorie des schémas \inserted{n'}apporte aucune simplification notable à la théorie \deleted{exposée} de B\deleted{orel}\handwritten{OREL}\deleted{\unsure{, tel qu'il|, a-t-il}} telle qu'elle est exposée dans BIBLE\handwritten{,} \deleted{C}\handwritten{c}'est pourquoi il n'a pas semblé utile de la reproduire \deleted{ici, notre objet étant} dans le présent Séminaire, notre but étant ici de déduire, de résultats connus sur les corps algébriquement clos, des résultats analogues valables sur tout préschéma de base. (Il\unsure{\,\textperiodcentered\,| }n'en sera plus de même pour la théorie de structure de C\deleted{hevalley}\handwritten{HEVALLEY}\deleted{\unsure{,}} des groupes semi-simples, qui\handwritten{,} semble-t-il\handwritten{,} se traite avec avantage ab ovo sur un préschéma de base quelconque).

\handwritten{Dans les exposés oraux \inserted{(que nous avons suivi dans les N\textsuperscript{os} 1 à 4)}, nous nous étions limités aux préschémas en groupes affines sur $S$, en nous appuyant de façon essentielle sur les théorèmes de représentabilité de Exp $\overline{\underline{\mathrm{XI}}}$ N\textsuperscript{o}4. Dans les présentes notes\deleted{\unsure{,}} (cf N\textsuperscript{os} 6 à 8) nous montrons rapidement comment on peut éliminer les hypothèses affines \deleted{\unsure{par une méthode|par une voie} \unsure{\illegible{2}|Exp XI}} \inserted{par une \deleted{\unsure{autre}} méthode} plus simple n'utilisant pas les résultats les plus délicats de Exp. $\overline{\underline{\mathrm{XI}}}$.}

% margin note written sideways
\inserted{\handwritten{\deleted{\unsure{Pour}}Pour d'autres généralisations des résultats contenus dans le présent exposé, voir également exposés \deleted{Exp.} $\overline{\underline{\mathrm{XV}}}$ et $\overline{\underline{\mathrm{XVI}}}$. (Il est évident que le contenu des exposés $\overline{\underline{\mathrm{XI}}}$, $\overline{\underline{\mathrm{XII}}}$, $\overline{\underline{\mathrm{XV}}}$, $\overline{\underline{\mathrm{XVI}}}$ devrait être complètement refondu).}}

% archive sheet number: 7
